\begin{lemma}{Sortierungsnorm}{Sortierungsnorm}
Seien $n \in \N$, $m \in \haken n$, 
\[\norm\cdot_{(m)}: \R^n \rightarrow [0, \infty[,
 v\mapsto\frac 1m\sum_{i=1}^m \mu_{i, n}\left(\pfeil {\abs\cdot}n(v) \right)\]
und
\[V\coloneqq  \mengea{z \in \R^n\setminus\meng 0}
    {\left(\forall\, i \in \haken n: \abs{z_i} \in\meng{0,\nicefrac 1m} \right)
    \wedge \left(\abs{\menge{i \in \haken n}{\abs{z_i} \neq 0}}\le m \right)}
    \text.\]
Dann ist $V$ endlich, nichtleer und es gilt:
\begin{enumerate}[(1)]
	\item \label{Sortierungsnorm1}
	$\sup \menge{\norm v_{2, n}}{v \in V} = \frac 1 {\sqrt m}$.
	\item \label{Sortierungsnorm2}
	$\forall\, w \in \R^n: \sup\menge{\sprod w v_n}{v \in V} = \norm w _{(m)}$.
	\item \label{Sortierungsnorm3}
	$\norm\cdot_{(m)}$ ist eine Norm auf dem $\R^n$ mit $\norm\cdot_{(m)} \le
	\norm\cdot_{\infty, n}$.
\end{enumerate}
\end{lemma}
\begin{proof}
Wegen $\frac 1 m e_1^n \in V \subseteq \meng{-\frac 1 m, 0, \frac 1 m}^n$ ist
$V$ nichtleer und endlich.\\
(1) Sei $v \in V$. Dann gilt:
\begin{align*}
\norm v_{2, n}^2 &= \sum_{i \in \haken n}\abs{v_i}^2
= \sum_{\substack{i\in \haken n\\ \abs{v_i} \neq 0}} \abs{v_i}^2
= \frac 1 {m^2} \sum_{\substack{i\in \haken n\\ \abs{v_i} \neq 0}} 1
\le \frac 1 {m^2} m = \frac 1 m\text.
\end{align*}
Also ist $\norm v _{2, n} \le\frac 1 {\sqrt m}$ und es gilt
\[v_0 \coloneqq  \sum_{i=1}^m \frac 1 m e_i^n \in V \quad \text{ und } \quad
\norm{v_0}_{2, n} = \frac 1 {\sqrt m}\text.\]
Somit gilt $\sup\menge{\norm w_{2, n}}{w \in V} = \frac 1 {\sqrt m}$.\\
(2) Sei $w \in \R^n$. Falls $w = 0$, so ist die Aussage trivial.\\
Sei also $w \neq 0$. Dann gilt für alle $v\in V$:
\begin{align*}\sprod w v_n& \le \sum_{i=1}^n \abs{w_i}\abs{v_i}
= \sum_{\substack{i \in \haken n\\ \abs{v_i} \neq 0} }\abs{w_i}\abs{v_i}
= \sum_{\substack{i \in \haken n\\ \abs{v_i} \neq 0} }\abs{w_i}\frac 1 m
\le \frac 1 m \sum_{i=1}^m \mu_{i, n}\left(\pfeil {\abs\cdot}n (w) \right)\\
&= \norm w_{(m)}\text.
\end{align*}
Ferner gibt es $\iota: \haken m \rightarrow \haken n$ injektiv mit
$\forall\, i \in \haken m: \abs{w_{\iota(i)}}
= \mu_{i, n}\left(\pfeil {\abs\cdot}n (w) \right)$.\\
Sei
\[v_w : \haken n \rightarrow \R,\, j \mapsto \begin{cases}
\frac 1 m \sgn(w_{j})&: j \in \Bild(\iota)\\
0 & : \text{ sonst}
\end{cases}.\]
Dann ist $v_w \in V$ und es gilt:
\begin{align*}
\sprod w {v_w}_n &= \sum_{i \in \haken n}w_i (v_w)_i
= \sum_{\substack{i \in \haken n\\ (v_w)_i \neq 0}}w_i (v_w)_i
= \sum_{i \in \Bild(\iota)}w_i \frac 1 m \sgn(w_i)\\
&= \frac 1 m\sum_{i \in \Bild(\iota)} \abs{w_i}
= \frac 1 m \sum_{i \in \haken m} \abs{w_{\iota(i)}}
= \frac 1 m \sum_{i=1}^m\mu_{i, n}\left(\pfeil {\abs\cdot}n (w) \right)
= \norm w _{(m)}\text.
\end{align*}
Somit ist auch (2) (analog zu (1)) gezeigt.\\
(3) Offenbar ist $\norm\cdot_{(m)}$ nichtnegativ--definit.\\
Seien $x, y \in \R^n$ und $t \in \R$. Dann gilt wegen $\abs t \ge 0$:
\begin{align*}
\norm{tx}_{(m)}&= 
\frac 1 m \sum_{i=1}^n\mu_{i, n}\left(\pfeil{\abs\cdot}n(tx) \right)
= \frac 1 m \sum_{i=1}^n\mu_{i, n}\left(\abs{t} \pfeil{\abs\cdot}n(x) \right)\\
&\overset{\refclaim{EigenschaftenDesKtenMaximums}
{EigenschaftenDesKtenMaximums3}}=
\frac 1 m\sum_{i=1}^n\abs{t}\mu_{i, n}\left(\pfeil{\abs\cdot}n(x) \right)
= \abs{t}\frac 1 m\sum_{i=1}^n\mu_{i, n}\left(\pfeil{\abs\cdot}n(x) \right)
= \abs{t}\norm{x}_{(m)}\text.
\end{align*}
Schließlich gilt:
\begin{align*}
\norm{x+y}_{(m)} &\overset{(2)}=
\sup_{v \in V} \sprod {x+y}v_n 
= \sup_{v \in V}\left(\sprod x v_n + \sprod y v_n \right)\\
&\le \left(\sup_{v \in V}\sprod x v_n \right) 
+ \left(\sup_{v \in V}\sprod y v_n \right)
\overset{(2)}= \norm x_{(m)} + \norm y _{(m)}\text.
\end{align*}
Also ist $\norm\cdot_{(m)}$ in der Tat eine Norm und es gilt
\[\norm\cdot _{(m)}
 = \frac 1 m \sum_{i=1}^m \mu_{i, n} \circ \pfeil{\abs\cdot}n
 \overset{\refclaim{EigenschaftenDesKtenMaximums}
 {EigenschaftenDesKtenMaximums1}}\le
 \frac 1 m \sum_{i=1}^m \mu_{1, n} \circ \pfeil{\abs\cdot}n
 = \mu_{1, n} \circ \pfeil{\abs\cdot}n = \norm\cdot_{\infty, n}\text.\]
\end{proof}